\section{Benevolent Centralized Power}

\begin{quotex}
These are three things which are rare indeed and are due to the grace of God -- namely, a human birth, the longing for Liberation, and the protecting care of a perfected sage.
\flright{\textsc{Shankara}}
\end{quotex}


I used to like to scan through the cable morning news shows now and then, not for the news, but to get the pulse of what we are supposed to be thinking. The shows are allegedly about the issues, but as we now know, it is never the issues, but rather affinity: whom we feel comfortable with. Fox and Friends is the more ``conservative'' of the shows. It is light and folksy. I do prefer it for the news because its only virtue is that it does demonstrate some \emph{token} opposition to the established powers. Nevertheless, it is not with them that I feel any affinity.\footnote{NOTE 2019: I no longer watch much of the network news.} The most liberal show is Morning Joe on MSNBC, so it is more interesting because the future can be found there.\footnote{NOTE 2021: It became unwatchable several years ago.} As the progressives think, and this show is the preparation for it, so goes the nation. Ironically, these are the people I would feel most comfortable with, if not for my self-imposed isolation. Having graduated from an elite New England university, I am quite familiar with that ``type''. They were the ones who dropped out the Logic class in the philosophy department, had trouble getting through Economics 101, and never even bothered with Science. Yet, as adults, they deem themselves competent to speak authoritatively on all those topics.

\textbf{Jacques Ellul} claimed that the purpose of education is to make it easier to indoctrinate people, so the highly educated are often the most malleable, or at least, they must be moulded first. So once the educated elite adopt a form of thought, the next layer below them follows suit and then on and on. The next great idea shows up on TV, movies, magazines, pop songs, and so on, until no one can remember any earlier time when the now accepted consensus opinion used to be taboo.

Needless to say, the cast members on Morning Joe are totally into control. They want to limit your choice of weapons, curtail your energy use, restrict what and how much you can eat, smoke, and so on, even how you think. I have felt, when listening to Mika, that they would like to treat the citizens the way Tyson raises chickens. That is, keep them in a cage to limit their motion, feed them a scientific diet, ignore their opinions and desires, and exterminate them when their useful lifetime has expired.

This week that suspicion was confirmed! When discussing the possibility of new sources of cheap energy for the masses, the multi-millionaire Donny Deutsch\footnote{\url{https://en.wikipedia.org/wiki/Donny_Deutsch}} claimed that we need ``\emph{benevolent centralized power}'' because anarchy will result from persons making their own decisions. Thus, this power can control our guns, what we eat and drink, our health care; they can tell us what to think and overturn what we previously had held to be dear. It is clear they have the power, and it is becoming more and more centralized. Benevolent, it is not.

\paragraph{The Element of Trust}


As we have already pointed out\footnote{\url{https://gornahoor.net/?p=295}}, an econo-socio-political system can be based on social relationships, liberal institutions, or on man himself. Consider this from \textbf{Valentin Tomberg}`s doctoral dissertation.

\begin{quotex}
A system, an order may be as useful and good as can be --- ultimately it depends on \textbf{people}, who abandon, relinquish, reinterpret, falsify, or simply betray it. And the number of a system's (loyal or disloyal) ``guardians'' need not be large. Often it takes only a few dozen people to destroy a generally recognized order. The big problem does not consist so much of which system to choose and how to organize the world, but of how and where to find \textbf{people}, who are truly of good will and who can be \textbf{trusted}. For ultimately, we are dealing with the element of trust. There is no system --- neither domestic nor international --- that can exist without trust as its final, practical, basic substance. For even with ten levels of control, there still must be --- as the \textbf{eleventh} level --- a \textbf{final} instance of control which can be trusted.
\end{quotex}


\paragraph{Uploading your soul}


There was some bad news\footnote{\url{http://io9.com/you-ll-probably-never-upload-your-mind-into-a-computer-474941498}} for those who are planning to upload a virtual copy of their brain into a computer where they can achieve some sort of immortality. I was a bit surprised to learn that the idea of a soul or spirit independent of matter is ``controversial'', but I don't travel in the smart circles. Nevertheless, most of those offering comments were not dissuaded. They are certain that their entire lives are actually computable processes and that whatever remains of subjective experience within them will be shown to be illusory. It is an intriguing prospect, since the quality of immortality will be independent of any moral consideration.

Yet I doubt their lives have gone as perfectly as they assume, but it may be easier to fix a ``bug'' in computer software than to actually change their own lives. But that is the problem. Whereas they seemed to have arisen spontaneously into existence, their virtual copy will require a programmer. We hope he turns out to be benevolent, or at least competent. I wonder if you will be able to choose your own virtual landscape when you are uploaded. \textbf{Dante} says it is we ourselves who choose our destiny, heaven or hell. I imagine the same notion will apply.

\paragraph{French Comments}


A couple of weeks ago, I got a long two thousand word comment on the Tarot Meditations blog\footnote{\url{https://wwww.meditationsonthetarot.com}}, all in French, from a women offering psychic readings. She was effusive in her praise, and asked questions about both metaphysics and blog software. Then she mentioned she hadn't bathed in three days. Apparently that must be an aphrodisiac to Frenchmen, but she neglected to take cultural diversity into account. So I rejected the comment. Some things perfume can't hide. So take it as a metaphor. Evil and perversity surround themselves with a sweet smell.

\paragraph{Shia and Iran}


The April 2013 issue of \emph{Culture Wars} has a long essay on E. Michael Jones' trip through Iran. He became a real Iranophile and noted something I myself mentioned here. He sees a certain affinity between Shia Islam and his own faith and brings up some useful points of similarity. The problem with Iran, he says, is the contraceptive mentality which had reduced the birthrate below replacement level and will open the door to feminism.

\paragraph{Library of Alexandria}


The saying goes that ``either what is in the library of Alexandria contradicts the Koran, in which case it is heretical and should be burnt, or it repeats what is in the Koran, in which case it is superfluous and should be burnt as well.'' That is sometimes the feeling I get from reading \textbf{Rene Guenon}. He is so focused on the ultimate liberation, that he seldom engages with more profane ideas (except negatively, as in his early books on Theosophy and Spiritism).

\textbf{Julius Evola}, on the other hand, engages the world (``the way of action''), incorporating the ideas of diverse figures, which need not be named here. That is why he appeals to more men today, or perhaps a different sort of man.

Coincidentally, I was reading about \textbf{Mouni Sadhu}'s experiences in the Indian ashram of \textbf{Sri Aurobindo}. He stumbled into the library and noticed all the books of Indian and European writers, most of which he was familiar with. He wrote:

\begin{quotex}
But now all these fascinating things had lost their charm for me. I realized that I was no longer interested in anything unconnected with my Path. It seemed as if knowledge of the Direct Path as shown by my Master subconsciously excluded all else. It meant that the desires of the mind, which is always eager to investigate everything, had begun to disappear.
\end{quotex}


At that point in his life, he realized he was no longer interested in them, as he was more focused on his own spiritual realization. I suppose there is a certain intellectual gluttony that desires to take in idea after idea. No doubt, for a certain type of man, there can be quite an intense ecstatic feeling from grasping some newly discovered thought. On the other hand, the man who had never had any such interest in learning previously in his life, is hardly on the same level.

\paragraph{Bhagavan Das}


Tomorrow I will be posting Evola's book review on the works of \textbf{Bhagavan Das}; to avoid clutter, I will put the introduction here. BD was an Indian Theosophist writing in the first half of the twentieth century. In true theosophical fashion, he claimed to have uncovered the Pranava Vada by Gargayana, which was dictated to him by a blind man who had memorized the million lines. Nevertheless, Evola has extraordinary high praise for the man and his works, claiming it was written for the Western mind. That is no doubt a side effect of his theosophical training, but he also relates Vedic ideas to Western thought, especially the German idealists.

This brings up a question, then. If the German idealists, who are actually the heirs of the Western tradition, knew so much or were so close to the same understanding, then why do we still bother to look to the East for enlightenment? Why didn't Evola himself promote his own works on Idealism as a spiritual path? Maybe Mouni Sadhu is correct, enough is enough. Or else, do you really need to hear the same thing over and over again?


\flrightit{Posted on 2013-04-27 by Cologero}

\begin{center}* * *\end{center}

\begin{footnotesize}
\begin{sffamily}
\texttt{Matt on 2013-04-27 at 11:06 said:}

Iran is not the only Islamic country with a birthrate that is well below the replacement level. Turkey is facing this issue as well, as are other countries in that region (it will have an even greater impact on these societies than European countries as at least European ones have the social structures to soften the impact of an aging population). Bernard Goldberg had a rather interesting take on this and what its effects are for relations between the modern west and the Islamic world.

And yes, a strong case can be made for the affinity between Shia Islam and Catholic/Orthodox Christianity. Makes one wonder if there would be less social tension in Europe if most muslim immigrants (specifically the young ones) were Shia instead of Sunni.

\hfill

\texttt{Avery Morrow on 2013-04-27 at 11:28 said:}

David P. Goldman has also written about this demographic issue at fascinating length. He said that in 300 years, Sub-Saharan Africans will be colonizing the dying embers of Europe in the name of equality and anti-racism, riding on rafts made from empty soda pop bottles, and all our metaphysical concerns are just determining the size and plot of our final descendants' graves.

Just kidding. He said that America will continue its healthy population growth indefinitely due to its God-granted mission.

\hfill

\texttt{Matt on 2013-04-27 at 12:17 said:}

Ah yes, it was Goldman (got him confused with Goldberg, as they both follow the classical liberal viewpoint).

\hfill

\texttt{Jason-Adam on 2013-04-27 at 16:13 said:}

As I see it, German Idealism is passive understanding at best, and does not offer a pathway to ``know'' the Ideas truly.

\hfill

\texttt{Jason-Adam on 2013-04-27 at 16:14 said:}

David Goldman is an ex-Larouchite

\hfill

\texttt{John on 2013-04-28 at 00:47 said:}

There's a sizable amount of so called trans-humanists with reactionary and counter-revolutionary political leanings on Twitter.

Such reasoning seems to me to be contradictory and absurd since they do not acknowledge a higher reality why bother with counter-revolution?

Trans-humanists are nothing more than what those who wrote the Vedas called ``demons who roam samsara denying the immortality of the soul'' since they seek to live forever in this unreal, reality.

\hfill

\texttt{Avery Morrow on 2013-04-28 at 11:52 said:}

One of Augustine's greatest works is a long list of mistakes he made in life.

\end{sffamily}
\end{footnotesize}

